(1) 对称轴角速度平方的横向部分为\(\dot\theta^2+\dot\phi^2\sin^2\theta\)，自旋分量为\(\dot\psi+\dot\phi\cos\theta\)，因此\(L=\frac I2(\dot\theta^2+\dot\phi^2\sin^2\theta)+\frac{I_3}2(\dot\psi+\dot\phi\cos\theta)^2-mgl\cos\theta\)。循环角\(\psi,\phi\)的正则动量分别为\(\lambda=p_\psi=I_3(\dot\psi+\dot\phi\cos\theta)\)与\(\mu=p_\phi=I\dot\phi\sin^2\theta+\lambda\cos\theta\)，均因Euler–Lagrange方程守恒。
(2) 从上述动量消去\(\dot\phi=(\mu-\lambda\cos\theta)/(I\sin^2\theta)\)与自旋速度\(\lambda/I_3\)，将常数自旋能\(\lambda^2/(2I_3)\)从总能减掉，得\(E'=\frac I2\dot\theta^2+U_{\mathrm{eff}}(\theta)\)，\(U_{\mathrm{eff}}=(\mu-\lambda\cos\theta)^2/(2I\sin^2\theta)+mgl\cos\theta\)。允许的\(\theta\)区间须满足\(E'\ge U_{\mathrm{eff}}\)；端点决定章动幅度。
(3) 规则进动有\(\dot\theta=0\)、\(\dot\phi=\Omega_p\)常量。倾角Lagrange方程给\(I\Omega_p^2\sin\theta\cos\theta-\lambda\Omega_p\sin\theta+mgl\sin\theta=0\)，故非竖直姿态满足\(I\cos\theta\,\Omega_p^2-\lambda\Omega_p+mgl=0\)。对\(\cos\theta\ne0\)，\(\Omega_p=[\lambda\pm\sqrt{\lambda^2-4Imgl\cos\theta}]/(2I\cos\theta)\)；判别式非负是给定倾角的存在条件。\(\cos\theta=0\)时直接由线性方程得到\(\Omega_p=mgl/\lambda\)，不应对分母为零的公式硬取值。
(4) 直立时需\(\mu=\lambda\)以免有效势在\(\theta=0\)发散。小角展开\(\cos\theta=1-\theta^2/2+\cdots\)、\((1-\cos\theta)^2/\sin^2\theta=\theta^2/4+O(\theta^4)\)，得到\(U_{\mathrm{eff}}=mgl+[\lambda^2/(8I)-mgl/2]\theta^2+O(\theta^4)\)。因此直立位置线性稳定须\(\lambda^2>4Imgl\)；等号时二次恢复刚度为零，必须作高阶判别。若\(g\to0\)，任意非零自旋满足该条件，符合无重力下的自由对称陀螺稳定性。
